\chapterimage{figs/chapterhead.png}
\chapter{GUI, Simulang, execução e observação do modelo}
\label{chap:gui_simulang_execucao_observacao}

\section{Introdução}
\label{sec:cap5_introducao}

Nos capítulos anteriores, o leitor colocou o \textit{GenESyS} em funcionamento, conheceu a lógica geral da interface gráfica e começou a trabalhar com modelos de exemplo já salvos no projeto. A partir deste ponto, torna-se possível dar um passo importante: sair da leitura apenas estrutural dos modelos e passar a observá-los também em execução, acompanhando sua dinâmica e relacionando a representação gráfica à representação textual do sistema.

Este capítulo tem exatamente esse propósito. Ele aprofunda o uso da GUI em três direções complementares. Primeiro, mostra como a execução do modelo deve ser acompanhada pelo usuário, não como uma operação opaca, mas como um processo observável. Segundo, introduz a aba \textit{Simulang} como forma textual de representação do mesmo modelo que já foi visto na interface gráfica. Terceiro, orienta a observação dos primeiros resultados, traços e sinais de comportamento do sistema, de modo que o usuário comece a construir uma leitura mais madura da simulação.

A importância dessa etapa é grande. Um simulador realmente útil não é apenas capaz de abrir modelos e executá-los, mas de permitir que o usuário compreenda o que está acontecendo. A GUI do \textit{GenESyS} foi pensada para oferecer diferentes perspectivas do mesmo trabalho: uma perspectiva gráfica, uma perspectiva textual e uma perspectiva dinâmica, ligada à execução. A partir daqui, o leitor começará a articular essas três camadas.

\section{Da estrutura estática à dinâmica do modelo}
\label{sec:cap5_estrutura_estatica_dinamica}

Quando o usuário abre um modelo e observa sua estrutura gráfica, ele vê uma descrição estática do sistema. Os componentes, conexões, dados de apoio e elementos gráficos indicam como a simulação foi organizada, mas ainda não mostram, por si sós, o comportamento temporal efetivo do modelo. A execução da simulação transforma essa estrutura estática em processo dinâmico.

Esse ponto é central para a compreensão do simulador. Um modelo não é apenas um diagrama bem organizado; ele é uma máquina de comportamento. A função do simulador é precisamente ativar essa máquina, conduzindo o avanço do tempo simulado, o processamento de eventos, a circulação de entidades e a atualização das informações relevantes do sistema.

Para o usuário, a principal consequência é esta: a execução não deve ser vista apenas como “apertar o botão de iniciar”. Ela deve ser entendida como a passagem da descrição do modelo para sua dinâmica operacional.

\subsection{O que muda quando o modelo começa a executar}
\label{subsec:cap5_o_que_muda_quando_executa}

Quando a simulação é iniciada, o usuário passa a trabalhar com um novo tipo de leitura. Antes, a pergunta principal era:
\begin{quote}
    Como este modelo está organizado?
\end{quote}

Durante a execução, essa pergunta muda para:
\begin{quote}
    Como este modelo se comporta ao longo do tempo?
\end{quote}

Isso significa observar:
\begin{itemize}
    \item se entidades estão sendo criadas, transformadas ou removidas;
    \item se o fluxo realmente percorre os elementos que o modelo sugere;
    \item se recursos, filas ou dados de apoio parecem ser utilizados de modo coerente;
    \item se o comportamento global da simulação faz sentido em relação à estrutura previamente lida.
\end{itemize}

\section{Comandos de execução e sua lógica de uso}
\label{sec:cap5_comandos_execucao_logica}

A GUI do \textit{GenESyS} disponibiliza um conjunto de comandos diretamente ligados à simulação. Entre eles, destacam-se os comandos de iniciar, pausar, retomar, avançar passo a passo, parar e configurar. Para o usuário, esses comandos não devem ser percebidos apenas como botões operacionais, mas como maneiras diferentes de se relacionar com a dinâmica do modelo.

\subsection{Iniciar a simulação}
\label{subsec:cap5_iniciar_simulacao}

O comando de início da simulação é o ponto de transição entre modelo estático e comportamento observado. Antes de acioná-lo, o usuário deve garantir que:
\begin{itemize}
    \item o modelo correto está carregado;
    \item a estrutura visual faz sentido;
    \item o sistema aparenta estar em condição de executar;
    \item o objetivo da observação está minimamente claro.
\end{itemize}

Essa última observação é importante. É melhor executar a simulação com alguma hipótese de leitura — por exemplo, observar um fluxo específico ou o efeito de certa propriedade — do que simplesmente iniciar a execução sem saber o que se pretende acompanhar.

\subsection{Pausar, retomar e avançar passo a passo}
\label{subsec:cap5_pausar_retormar_passo_a_passo}

Uma das grandes vantagens de uma GUI rica é permitir diferentes ritmos de observação da simulação. Em algumas situações, o usuário quer apenas verificar o comportamento global. Em outras, deseja observar trechos específicos do fluxo com mais atenção. É nesse contexto que comandos como pausar, retomar e executar passo a passo se tornam particularmente valiosos.

Esses comandos transformam a simulação em objeto de inspeção. Em vez de apenas deixar o sistema correr até o fim, o usuário pode:
\begin{itemize}
    \item interromper temporariamente a execução;
    \item analisar o estado atual;
    \item retomar a simulação;
    \item ou avançar de forma controlada.
\end{itemize}

Esse modo de uso é especialmente útil em modelos de aprendizagem, em depuração inicial e na leitura do comportamento de elementos cuja função ainda não esteja totalmente clara para o usuário.

\subsection{Parar e configurar}
\label{subsec:cap5_parar_configurar}

Os comandos de parada e configuração também merecem atenção. Parar a simulação não é apenas encerrar um processo, mas interromper a evolução do modelo de forma controlada. Já configurar a simulação significa ajustar condições de execução antes do início do experimento.

Embora o aprofundamento desses parâmetros dependa do modelo e do contexto, o usuário deve reter desde já que a execução não é completamente rígida: ela pode ser preparada, iniciada, observada e interrompida segundo necessidades diferentes.

\section{Observando o comportamento do modelo durante a execução}
\label{sec:cap5_observando_comportamento_execucao}

A observação da execução deve ser conduzida de forma orientada. Se o usuário tentar acompanhar simultaneamente todos os sinais da GUI, toda a leitura se tornará difusa. O ideal é observar o modelo em níveis.

\subsection{Primeiro nível: o fluxo global}
\label{subsec:cap5_primeiro_nivel_fluxo_global}

No primeiro nível, o usuário deve observar o comportamento do modelo de forma global. As perguntas mais importantes são:
\begin{itemize}
    \item a simulação iniciou corretamente?
    \item o fluxo principal parece coerente com a estrutura do modelo?
    \item o sistema evolui de forma plausível?
    \item o término da execução faz sentido?
\end{itemize}

Essa é uma leitura macroscópica. Ela não busca detalhes finos, mas consistência geral.

\subsection{Segundo nível: elementos específicos}
\label{subsec:cap5_segundo_nivel_elementos_especificos}

Depois de observar o comportamento global, o usuário pode concentrar-se em elementos específicos do modelo. Isso pode incluir:
\begin{itemize}
    \item um componente cuja função ainda não esteja clara;
    \item um ponto do fluxo onde o comportamento parece particularmente importante;
    \item uma fila, recurso ou dado de apoio relevante;
    \item uma propriedade cujo efeito foi recentemente modificado.
\end{itemize}

Essa leitura localizada é mais produtiva depois que a dinâmica geral já foi compreendida minimamente.

\subsection{Terceiro nível: comparação entre execuções}
\label{subsec:cap5_terceiro_nivel_comparacao_execucoes}

Uma prática especialmente útil é comparar a execução do modelo antes e depois de pequenas alterações. Esse procedimento ajuda o usuário a compreender não apenas como o sistema funciona, mas também como reagirá a mudanças. A GUI, nesse sentido, não é apenas interface de observação, mas também laboratório de comparação.

\section{A aba \textit{Simulang} como representação textual do modelo}
\label{sec:cap5_aba_simulang}

Depois de o leitor já ter visto o modelo na GUI e já tê-lo observado em execução, a representação textual passa a fazer muito mais sentido. É exatamente por isso que a aba \textit{Simulang} aparece agora, e não antes.

A função da aba \textit{Simulang} não é substituir a GUI, mas oferecer uma segunda forma de leitura do mesmo modelo. Nessa forma de leitura, o usuário deixa de ver apenas a organização gráfica do sistema e passa a enxergar uma descrição textual correspondente. Isso é valioso por várias razões:

\begin{itemize}
    \item torna o modelo mais explícito em termos formais;
    \item permite comparar a representação visual com a representação textual;
    \item ajuda a consolidar a compreensão do modelo;
    \item prepara o caminho para capítulos posteriores sobre linguagem e desenvolvimento.
\end{itemize}

\subsection{Como ler a aba \textit{Simulang}}
\label{subsec:cap5_como_ler_simulang}

A leitura da aba \textit{Simulang} deve ser feita com a mesma cautela didática usada no restante da primeira parte do livro. O usuário não precisa, neste momento, dominar toda a sintaxe ou toda a semântica da linguagem. O que ele precisa é perceber que:
\begin{itemize}
    \item a descrição textual corresponde ao modelo que já viu na GUI;
    \item elementos visuais do modelo possuem uma expressão textual;
    \item a linguagem oferece uma visão mais formal do sistema modelado.
\end{itemize}

A forma mais produtiva de iniciar essa leitura é comparar trechos da representação textual com elementos já conhecidos da interface gráfica. Em vez de tentar decifrar o texto inteiro de uma vez, o usuário deve perguntar:
\begin{quote}
    Que parte do modelo gráfico esta região do texto parece descrever?
\end{quote}

\subsection{Por que a representação textual é útil ao usuário}
\label{subsec:cap5_por_que_simulang_util_ao_usuario}

Mesmo que o usuário continue trabalhando prioritariamente pela GUI, a representação textual lhe oferece benefícios concretos. Ela ajuda a:
\begin{itemize}
    \item consolidar a compreensão do modelo;
    \item tornar mais explícita a estrutura do sistema;
    \item perceber o modelo além da sua disposição visual;
    \item começar a pensar na relação entre interface e linguagem.
\end{itemize}

Essa experiência é particularmente valiosa porque mostra, desde cedo, que o \textit{GenESyS} não trata o modelo apenas como desenho, mas também como descrição formal.

\section{Tracing, breakpoints e observabilidade inicial}
\label{sec:cap5_tracing_breakpoints_observabilidade}

Uma característica importante de um simulador bem estruturado é oferecer mecanismos de observação do comportamento do sistema em execução. No \textit{GenESyS}, isso se expressa por diferentes meios disponíveis na GUI, incluindo traços, eventos, relatórios e mecanismos de parada ou acompanhamento controlado.

Nesta etapa do manual, o leitor não precisa ainda dominar profundamente esses mecanismos. Mas precisa entender sua função geral: eles tornam a simulação mais observável.

\subsection{Tracing como leitura do comportamento}
\label{subsec:cap5_tracing_leitura_comportamento}

O tracing deve ser visto como uma forma de registro do que ocorre na simulação. Ele ajuda o usuário a acompanhar o modelo não apenas pela percepção visual da interface, mas também por uma narrativa operacional do que está acontecendo.

Isso é especialmente útil quando o comportamento do modelo não é imediatamente evidente apenas pela área gráfica. Em tais casos, o tracing ajuda a responder:
\begin{itemize}
    \item que eventos estão ocorrendo;
    \item em que ordem;
    \item com que efeitos aparentes sobre o estado do sistema.
\end{itemize}

\subsection{Breakpoints como ferramenta de observação controlada}
\label{subsec:cap5_breakpoints_observacao_controlada}

Os breakpoints, quando utilizados, permitem interromper a simulação em pontos de interesse. Para o usuário, isso transforma a execução em objeto examinável. Em vez de deixar o modelo evoluir completamente sem intervenção, pode-se parar em estados relevantes e observar melhor o que está acontecendo.

Essa ferramenta é particularmente útil quando:
\begin{itemize}
    \item se deseja estudar o comportamento de um trecho específico do modelo;
    \item se quer validar o efeito de uma modificação;
    \item se pretende compreender melhor a dinâmica de um elemento específico.
\end{itemize}

\section{Primeira leitura dos resultados}
\label{sec:cap5_primeira_leitura_resultados}

Além de observar a execução em andamento, o usuário precisa começar a desenvolver o hábito de interpretar os resultados produzidos pelo simulador. Aqui, novamente, a orientação deve ser progressiva.

A primeira leitura de resultados não deve ser estatisticamente sofisticada nem excessivamente ambiciosa. Ela deve responder a perguntas simples:
\begin{itemize}
    \item a simulação se comportou como esperado?
    \item houve sinais coerentes de progressão e término?
    \item mudanças simples realizadas no modelo produziram efeitos observáveis?
    \item os relatórios ou sinais apresentados pela GUI parecem compatíveis com a estrutura do modelo?
\end{itemize}

Só depois dessa etapa inicial é que se torna útil avançar para leituras mais técnicas.

\subsection{Resultados como continuação da leitura do modelo}
\label{subsec:cap5_resultados_continuacao_leitura}

Um ponto importante é que os resultados não devem ser tratados como algo separado do modelo. Eles são continuação da leitura do modelo por outros meios. O modelo gráfico mostra estrutura, a execução mostra dinâmica, e os resultados mostram consequências observáveis dessa dinâmica. Essas três dimensões devem ser sempre articuladas.

\section{Boas práticas para observar modelos em execução}
\label{sec:cap5_boas_praticas_observacao_execucao}

Algumas práticas tornam a observação da simulação muito mais proveitosa.

\subsection{Defina um foco de observação}
\label{subsec:cap5_defina_foco_observacao}

Antes de iniciar a execução, procure decidir o que observar primeiro: o fluxo global, um componente específico, um efeito de certa alteração, ou a relação entre GUI e Simulang. Isso evita dispersão.

\subsection{Alterne entre visões}
\label{subsec:cap5_alterne_entre_visoes}

A força da GUI do \textit{GenESyS} está em oferecer múltiplas perspectivas sobre o mesmo modelo. Alterne entre:
\begin{itemize}
    \item estrutura gráfica;
    \item execução observada;
    \item traços ou eventos;
    \item representação textual em \textit{Simulang}.
\end{itemize}

Essa alternância costuma produzir compreensão mais profunda do que qualquer visão isolada.

\subsection{Faça pequenas modificações e compare}
\label{subsec:cap5_pequenas_modificacoes_compare}

Sempre que possível, faça pequenas alterações controladas no modelo e compare o comportamento antes e depois. Essa estratégia transforma a observação em experimento guiado.

\subsection{Não tente esgotar tudo em uma única execução}
\label{subsec:cap5_nao_esgotar_tudo}

Uma simulação raramente revela tudo o que o modelo pode ensinar em uma única execução. É melhor fazer leituras sucessivas, cada uma com foco mais claro, do que tentar absorver tudo ao mesmo tempo.

\begin{table}[htbp]
    \centering
    \caption{Estratégia recomendada para observar modelos em execução no GenESyS.}
    \label{tab:cap5_estrategia_observacao_execucao}
    \small
    \begin{tabular}{|p{4.3cm}|p{6.9cm}|}
        \hline
        \textbf{Etapa} & \textbf{Objetivo} \\
        \hline
        Observar a estrutura gráfica & Compreender a organização estática do modelo \\
        \hline
        Iniciar a simulação & Transformar a estrutura em comportamento dinâmico \\
        \hline
        Acompanhar o fluxo global & Verificar consistência geral da execução \\
        \hline
        Explorar tracing e eventos & Tornar o comportamento mais observável \\
        \hline
        Ler a aba \textit{Simulang} & Relacionar o modelo gráfico à sua descrição textual \\
        \hline
        Comparar execuções & Entender o efeito de pequenas alterações no modelo \\
        \hline
    \end{tabular}
\end{table}

\section{Considerações Finais}
\label{sec:cap5_consideracoes_finais}

Este capítulo aprofundou o uso da GUI do \textit{GenESyS} em uma direção essencial: a articulação entre estrutura visual do modelo, execução da simulação, observabilidade do comportamento e representação textual em \textit{Simulang}. A principal ideia que o leitor deve reter é que essas dimensões não competem entre si. Elas se complementam. O modelo gráfico permite ver a organização do sistema, a simulação permite observar sua dinâmica, os mecanismos de tracing e pausa tornam esse comportamento mais legível, e a aba \textit{Simulang} oferece uma forma textual de consolidar essa compreensão.

Com isso, a primeira parte do manual começa a ganhar densidade real de uso. O leitor já não apenas abriu a GUI e carregou um modelo; ele passou a observar o simulador em funcionamento e a relacionar suas várias representações. O próximo passo natural, a partir daqui, é consolidar o uso do sistema com exemplos mais orientados e preparar a transição para uma compreensão mais madura da relação entre modelagem, simulação e estrutura interna do GenESyS.

Teste seu entendimento desse conteúdo respondendo as seguintes perguntas:

\begin{enumerate}
    \item Por que a execução de um modelo deve ser entendida como passagem da estrutura estática para a dinâmica do sistema?

    \item Qual é a utilidade de comandos como pausar, retomar e avançar passo a passo durante a simulação?

    \item Por que a aba \textit{Simulang} só deve ser introduzida depois que o usuário já viu e executou o modelo na GUI?

    \item Como tracing, breakpoints e resultados complementam a observação visual da execução do modelo?
\end{enumerate}